\chapter*{Introduction générale}
\addcontentsline{toc}{chapter}{Introduction générale}

\section*{Contexte général}
\addcontentsline{toc}{section}{Contexte général}

Dans le domaine de la santé, l'accès à une information fiable peut avoir un impact direct sur le service
rendu au citoyen. Pour les pharmacies, cette information concerne notamment l'adresse, le numéro de
téléphone, les horaires et les périodes de garde. Lorsqu'un utilisateur cherche une pharmacie disponible,
il attend une réponse simple et rapide, surtout en dehors des horaires habituels.

En Tunisie, les informations relatives aux pharmacies de garde existent, mais leur consultation reste
parfois peu pratique. Elles peuvent être diffusées sous forme de listes, relayées localement ou mises à
jour manuellement. Cette dispersion complique la recherche depuis un téléphone, en particulier lorsque
l'utilisateur ne connaît pas bien la ville ou le gouvernorat concerné.

Le projet \textbf{PharmacieConnect} s'inscrit dans ce contexte. Il associe une API backend FastAPI, un
tableau de bord administrateur React/Vite, une application mobile Expo/React Native et un assistant
conversationnel spécialisé en premiers secours. Le backend porte la logique commune, tandis que les
interfaces répondent à des usages complémentaires : la consultation citoyenne, l'administration des
données et l'assistance médicale d'urgence. L'assistant conversationnel enrichit l'application mobile en
fournissant des réponses rapides et fiables aux questions de premiers secours, complément naturel au
repérage des pharmacies de garde.

\section*{Problématique}
\addcontentsline{toc}{section}{Problématique}

La problématique du projet peut être formulée ainsi : comment concevoir une plateforme qui permette au
citoyen de trouver rapidement une pharmacie ou une pharmacie de garde, tout en offrant à
l'administrateur un moyen fiable de gérer, corriger et suivre les données ?

Cette question impose plusieurs contraintes. L'application mobile doit rester simple à utiliser et fournir
des résultats lisibles. Le tableau de bord doit faciliter l'import, la correction et le suivi des données.
Le backend, quant à lui, doit contrôler les accès, valider les informations reçues et conserver une trace
des actions importantes.

\section*{Objectifs du projet}
\addcontentsline{toc}{section}{Objectifs du projet}

Le projet vise à réaliser une plateforme complète autour des pharmacies et des pharmacies de garde, enrichie par un assistant de premiers secours. Les objectifs retenus sont les suivants :

\begin{itemize}
    \item concevoir une API backend centralisant les données et les règles métier ;
    \item mettre en place une authentification sécurisée avec des jetons \gls{jwt} ;
    \item mettre à disposition une interface web pour gérer les pharmacies, les gardes, les médicaments et les utilisateurs autorisés ;
    \item réaliser une application mobile orientée vers la recherche de pharmacies, la consultation des gardes et l'accès au catalogue de médicaments ;
    \item intégrer un assistant conversationnel spécialisé en premiers secours utilisant la récupération d'informations hybride ;
    \item intégrer une recherche par proximité à partir des coordonnées \gls{gps} ;
    \item prendre en charge l'import de fichiers \gls{csv} avec validation des données et retour d'erreurs ;
    \item ajouter des journaux d'audit, des indicateurs et un contrôle d'accès basé sur les rôles ;
    \item tester les fonctionnalités principales, y compris l'assistant et ses mécanismes de sécurité, à travers des tests backend et des scénarios fonctionnels.
\end{itemize}

\section*{Méthodologie adoptée}
\addcontentsline{toc}{section}{Méthodologie adoptée}

La démarche suivie est itérative et incrémentale, avec une organisation inspirée de Scrum
\cite{scrumguide}. Le travail a été découpé en cycles courts afin d'avancer progressivement sur des
modules vérifiables : analyse du besoin, étude de l'existant, définition des fonctionnalités, conception de
l'architecture, développement du backend, développement de l'administration web, développement de
l'application mobile, intégration, tests et rédaction du rapport.

La modélisation s'appuie sur le formalisme \gls{uml} \cite{uml} pour représenter les acteurs, les cas
d'utilisation, les classes principales et les échanges entre les composants. La réalisation a d'abord
porté sur les règles métier du backend, puisque l'API centralise la validation, la sécurité et l'accès aux
données. Les interfaces web et mobile ont ensuite été raccordées progressivement aux endpoints
disponibles.

Les fonctionnalités sensibles ont été vérifiées au fur et à mesure de leur intégration. Cette progression a
facilité la validation des imports, de l'authentification, de la gestion des rôles, de la recherche par
proximité et des principaux parcours utilisateur.

\section*{Organisation du rapport}
\addcontentsline{toc}{section}{Organisation du rapport}

Le rapport est organisé en quatre chapitres. Le premier situe le projet dans son cadre général, présente
le domaine pharmaceutique et analyse l'existant. Le deuxième précise les acteurs, les besoins
fonctionnels et non fonctionnels, y compris les besoins relatifs à l'assistant de premiers secours, et expose
les règles de gestion. Le troisième est consacré à la conception : architecture globale, architecture de
l'assistant, modèle de données, diagrammes \gls{uml} et sécurité. Le quatrième décrit la réalisation
technique du backend, du tableau de bord, de l'application mobile et de l'assistant, ainsi que
l'intégration, la configuration, les tests et les difficultés rencontrées. Une conclusion générale, une
bibliographie et des annexes complètent le document.
